% Source: https://artofproblemsolving.com/community/c3h1871927_cue_ball_geometry_problem
% Stable identity: aops-topic-1871927-q01
When the ball touches the rail, its centre is $1$ inch from the rail. So the centre moves in straight lines and reflects off the line $\ell$ parallel to the rail at distance $1$ (dashed below); $A$ and $B$ are $6-1=5$ and $9-1=8$ inches from $\ell$. The ball touches the rail at the point $K$ directly beyond the centre's turning point $T$ on $\ell$.

Measured parallel to the rail, $A$ and $B$ are $\sqrt{5^2-(9-6)^2}=4$ inches apart.

\begin{asy}
usepackage("amsmath");
size(7cm);
defaultpen(fontsize(10pt));
real t=20/13;
pair A=(0,6),B=(4,9),T=(t,1),K=(t,0);
fill((-1.2,-0.35)--(5.4,-0.35)--(5.4,0)--(-1.2,0)--cycle,gray(0.75));
draw((-1.2,0)--(5.4,0),linewidth(1.2));
draw((-1.2,1)--(5.4,1),dashed+gray(0.4));
draw(circle(T,1),heavygreen);dot(T,heavygreen);
draw(A--T--B,blue+linewidth(1));
draw(A--(0,-1.3),dotted);draw(B--(4,-1.3),dotted);draw(T--K,dotted);draw(K--(t,-0.45),dotted);
dot(A);dot(B);dot(K,red);
label("$A$",A,W);label("$B$",B,E);label("$T$",T+(0.5,0.42));label("$K$",(t,-0.72),red);
label("$5$",(0,3.5),W);label("$8$",(4,5),E);label("$1$",(0,0.5),W);label("rail",(5.4,-0.2),E);
draw((t,-1.3)--(4,-1.3),red,Arrows(6));label("$\dfrac{32}{13}$",((t+4)/2,-1.3),S,red);
draw((0,-1.3)--(t,-1.3),gray(0.4),Arrows(6));label("$\dfrac{20}{13}$",(t/2,-1.3),S,gray(0.4));
label("\mbox{}",(min(currentpicture,user=true)+max(currentpicture,user=true))/2,black); // keeps the text after the figure black when LaTeX typesets the labels inline
\end{asy}

Since the angles that $AT$ and $TB$ make with $\ell$ are equal, the right triangles formed by $A$, $T$ and the foot of $A$ on $\ell$, and by $B$, $T$ and the foot of $B$ on $\ell$, are similar, with legs perpendicular to $\ell$ equal to $5$ and $8$. So $T$ divides the horizontal distance $4$ in the ratio $5:8$:
\[
\text{horizontal distance from }B\text{ to }T=4\cdot\dfrac{8}{13}=\dfrac{32}{13}.
\]
The point of the rail nearest to $B$ is the foot of the perpendicular from $B$, and $K$ is directly below $T$. So the ball should hit the rail $\boxed{\dfrac{32}{13}\text{ inches}}\approx2.46$ inches from that point (and $\dfrac{20}{13}$ inches from the point nearest to $A$).
